
\section{图表示例}

在此输入本章的引言内容，概述本章的主要内容和目的。

\subsection{图片插入示例}

在此输入引出图片的过渡性文字。图 \ref{fig:example_fig} 
展示了相关示意图（此处使用黑色色块作为占位符）。

\begin{figure}[H]
    \centering
    % 占位符：实际使用时请替换为如下命令
    % \includegraphics[width=0.6\linewidth]{figures/your_image.png}
    \rule{0.6\linewidth}{5cm}
    \caption{图片标题}
    \label{fig:example_fig}
\end{figure}

在此输入对图片的分析和说明文字，解释图片所展示的
内容及其意义。

\subsection{表格插入示例}

在此输入引出表格的过渡性文字，说明表格的用途和内容。
如表 \ref{tab:chapter3_table} 所示，展示了相关数据。

\begin{table}[H]
    \centering
    \caption{示例表格标题}
    \label{tab:chapter3_table}
    \begin{tabularx}{0.95\textwidth}{l X X}
        \toprule
        列标题1 & 列标题2 & 列标题3 \\
        \midrule
        数据1 & 数据描述内容 & 备注说明 \\
        数据2 & 数据描述内容 & 备注说明 \\
        数据3 & 数据描述内容 & 备注说明 \\
        \bottomrule
    \end{tabularx}
\end{table}

\subsection{本章小结}

在此总结本章展示的图表使用方法和注意事项。
